\section{الفصل~\ref{sec:chemoutput}. المخرج الكيميائي}

السلكان إلى القلب واختلاف سرعتيهما، والتغذية الراجعة السالبة في سلاسل الهرمونات، والشفرة بتردد الدفعات، والمولّد اليومي وضبطه بالضوء، كلها مُقاسة مباشرة؛ والحد $K<8$ مبرهنة من نظرية التحكم مشتقة في الفصل؛ أما النبضات التي تولّدها التغذية الراجعة للسلسلة نفسها ففرضية تدعمها تجربة قوية.

{\footnotesize
\setlength{\tabcolsep}{3pt}\renewcommand{\arraystretch}{1.15}
\begin{longtable}{>{\raggedright}p{0.5cm}>{\raggedright}p{4.0cm}>{\raggedright}p{5.6cm}>{\raggedright}p{2.3cm}>{\raggedright\arraybackslash}p{1.7cm}}
\toprule
م & القضية & التجربة وما قيس & المصدر & الحالة \\
\midrule\endfirsthead
\toprule
م & القضية & التجربة وما قيس & المصدر & الحالة \\
\midrule\endhead
1 & ناظمة إيقاع القلب تنبض وحدها نحو $100$ مرة في الدقيقة؛ في الراحة يكون المكبح مضغوطًا، والتردد عادةً نحو $60\text{--}70$ & بشر، إيقاف السلكين إلى القلب إيقافًا تامًا: التردد الذاتي أعلى من تردد الراحة ويتناقص مع العمر، وهو عند البالغين نحو $100$ ضربة في الدقيقة. إذن يغلب المكبح في الراحة. تردد الراحة $60\text{--}70$ قيمة نموذجية لم يُستشهد لها بمصدر مستقل. & \cite{JoseCollison1970ihr} & مُقاس \\
2 & المسرِّع \nt{NORA} والمكبح \nt{ACET} مرشحا تمرير منخفض، والمكبح أسرع~\eqref{eq:gasbrake} & كلب، تحفيز معدَّل التردد للعصب المبهم أو للعصب الودي القلبي ودالة التحويل إلى تردد الأذين: كلا الطريقين مرشح تمرير منخفض، لكن تردد القطع في الطريق الودي أدنى بكثير وتضاف إليه تأخيرة صافية نحو $1.7\,\text{s}$. الخصائص تعتمد على متوسط التوتر، فالصيغة الخطية~\eqref{eq:gasbrake} تقريب. & \cite{Berger1989} & مُقاس \\
3 & المكبح سريع لأن \nt{ACET} يفتح قناة البوتاسيوم دون مرسال ثانٍ، أما \nt{NORA} فعبر سلسلة تفاعلات & رقع غشائية من خلايا الأذين: قناة البوتاسيوم التي يفتحها الأسيتيل كولين تفتحها الوحدتان الفرعيتان $\beta\gamma$ لبروتين \foreignlanguage{english}{G} المرتبط بالمستقبل، دون مرسال ثانٍ داخل الخلية. طريق \nt{NORA} عبر \foreignlanguage{english}{cAMP} لم يُستشهد له هنا. & \cite{Logothetis1987gbg} & مُقاس \\
4 & يدور الدم في الجسم في نحو دقيقة؛ ويصل \nt{ADRE} من مجرى الدم إلى كل الأعضاء التي تملك المستقبل & تقدير عام لرتبة المقدار؛ وهكذا صيغ في الفصل، ولم يُستشهد بمصدر. & --- & تقدير \\
5 & تحيط بسلسلة الهرمونات تغذية راجعة سالبة: الهرمون الأخير يثبّط المراحل الأولى & مراجعة للتجارب: هرمونات قشرة الكظر تثبّط إفراز الغدة النخامية لـ\foreignlanguage{english}{ACTH} وإفراز الوطاء للهرمون المطلِق في ثلاثة نطاقات زمنية: سريع ومتوسط وبطيء. & \cite{KellerWood1984fb} & مُقاس \\
6 & ثلاث مراحل متماثلة مستقرة عند $K<8$~\eqref{eq:cascadestable}، والدور عند الحد $2\pi\tau/\sqrt3\approx3.6\tau$ & مشتق في الفصل: عند التردد $\sqrt3/\tau$ تعطي المراحل الثلاث إزاحة طور $180^\circ$ وتوهينًا بـ$8$ أضعاف. & اشتقاق في الفصل & نظرية \\
7 & خطأ المستوى $1/(1+K)$، لذلك لا يثبّت هذا المنظِّم المستوى بدقة أفضل من نحو $10\,\%$ (القضية~\ref{prop:cascade}) & نتيجة للسطر~6 في التغذية الراجعة التناسبية. المحور الحقيقي فيه تغذية راجعة على عدة مراحل وفي عدة نطاقات زمنية (السطر~5)، فالحد يخص النموذج~\eqref{eq:cascade} ولم يُقَس. & اشتقاق في الفصل & نظرية \\
8 & عند $K=4$ و$7$ يستقر المستوى، وعند $K=12$ نبضات منتظمة بدور $3.7\text{--}3.9\,\tau$ (الشكل~\ref{fig:cascade}) & حساب بـ\texttt{models/\allowbreak chem\_\allowbreak models.py}: عند $K=12$ الأدوار المتتالية $3.9$ و$3.7$ و$3.8$ و$3.7\,\tau$؛ وعند $K=4$ السعة في نهاية الحساب $0.003$. & \texttt{models/\allowbreak chem\_\allowbreak models.py} & نموذج \\
9 & يخرج هرمون الإجهاد نبضاتٍ نحو مرة في الساعة؛ وتولّد النبضاتِ التغذيةُ الراجعة نفسها، دون مولّد مستقل & جرذان، تسريب مستمر للهرمون المطلِق: يتذبذب \foreignlanguage{english}{ACTH} والكورتيكوستيرون بتردد فيزيولوجي يقارب مرة في الساعة، فلا حاجة إلى مدخل إيقاعي من الوطاء للنبضات. ونموذج النخامية والكظر بتغذية راجعة متأخرة يعطي مثل هذه التذبذبات. & \cite{Walker2012glc, Walker2010} & فرضية \\
10 & المتلقّي يقرأ تردد الدفعات لا المستوى & قرود بلا إفراز ذاتي للهرمون المطلِق لموجّهات الغدد التناسلية: التسريب المستمر لا يعيد إفراز هرمونات النخامية، والدفعات مرة كل ساعة تعيده؛ والانتقال إلى التسريب المستمر يخمد الاستجابة من جديد. تعب المستقبل بوصفه مرشح تمرير عالٍ تفسير. & \cite{Belchetz1978} & مُقاس \\
11 & الترددات المختلفة للدفعات تؤثر تأثيرات مختلفة في خلايا مختلفة من الغدة نفسها & القرود نفسها: زيادة تواتر الدفعات إلى $2\text{--}5$ في الساعة تخفض هرموني النخامية كليهما؛ وتقليله إلى دفعة كل $3$ ساعات يخفض أحدهما (\foreignlanguage{english}{LH}) ويرفع الآخر (\foreignlanguage{english}{FSH})، فالنسبة بينهما يحددها التردد. & \cite{Wildt1981gnrh} & مُقاس \\
12 & يولّد الإيقاعَ اليومي تغذيةٌ راجعة سالبة داخل الخلية بتأخر عدة ساعات & مراجعة: بروتينات جينات الساعة تثبّط بتأخر نسخ جيناتها هي؛ والحلقة من النسخ والترجمة تعطي دورًا يقارب يومًا. & \cite{TakahashiJS2017} & مُقاس \\
13 & المولّد الرئيسي مجموعة من عشرات آلاف العصبونات تزامن بقية الجسم & مراجعة: النواة فوق التصالبة هي المولّد اليومي الرئيسي؛ عصبوناتها المفردة في المزرعة مولّدات مستقلة، والشبكة تزامنها؛ وعند الفأر نحو $10^4$ عصبون في كل جانب. & \cite{Welsh2010scn} & مُقاس \\
14 & الدور الذاتي عند الإنسان $24.18$ ساعة & بشر في ضوء خافت وبجدول منفصل عن اليوم: دور إيقاعات الميلاتونين والحرارة والكورتيزول $24.18$ ساعة في المتوسط عند الشباب والمسنين، بتشتت ضيق. & \cite{Czeisler1999} & مُقاس \\
15 & الضوء يضبط المولّد كحلقة مقفلة الطور~\eqref{eq:pll}؛ والإزاحة اللازمة نحو $11$ دقيقة في اليوم & بشر، ضوء ساطع مفرد مدته $6.7$ ساعة في أوقات مختلفة من اليوم: منحنى استجابة طورية من النوع الأول بسعة $5$ ساعات، تأخير قبل الحد الأدنى لحرارة الجسم وتقديم بعده. المعادلة~\eqref{eq:pll} نموذج معياري؛ و$11$ دقيقة $=0.18$ ساعة حساب بسيط. & \cite{Khalsa2003prc} & مُقاس \\
16 & بلا ضوء لا تزامن، وينجرف الإيقاع إلى دوره الذاتي & مكفوفون كليًا لا يدركون الضوء ويعيشون في المجتمع العادي: عند نحو نصفهم تسير إيقاعات الميلاتونين والكورتيزول بدورها الخاص الذي لا يطابق اليوم. & \cite{Sack1992blind} & مُقاس \\
\bottomrule
\end{longtable}}
